% Appendix G -- Strong Baselines (body). Sources: main/reports/E5_tdmpc2.md, E5_dinowm.md, posthoc_tdmpc2.md;
%   tmlr_ext/MT_317m/REPORT.md 5, DW_dinowm/REPORT.md 5, 9.2. Table: tables/tabG_h3b.tex.
\textbf{Post-training quantizers (E5).} E5 asked whether dither still helps on top of the best post-training quantizer at the same bit width. Each TD-MPC2 cell chose among four quantizers: per-channel ranges, MSE-optimal clipping, SmoothQuant-style smoothing, and a mixed-precision grid that keeps the last dynamics and reward layers at 8 bits; each DINO-WM cell among seven: per-channel ranges, percentile ranges at 99.9, 99.99 (the E1 grid) and 99.999, MSE-optimal clipping, smoothing, and the last block at 8 bits. The quantizer $B$ was chosen on held-out data by the smallest damage of $B$+\dtm (TD-MPC2, held-out seeds) or the highest success (DINO-WM, 20 held-out pairs), ties broken by list order; per-channel ranges won \nv{g.sel.td.pc} of the \nv{g.sel.td.n} TD-MPC2 cells. Table~\ref{tab:h3b} gives every cell. The best quantizer already removes most of \dtm's damage, so the room left for dither is small: the extra gain of \shiter has a lower bound above zero in \nv{e5.h3b.td} of \nv{e5.h3b.td.of} TD-MPC2 cells, exactly the threshold, and in \nv{e5.h3b.dw} of \nv{e5.h3b.dw.of} DINO-WM cells. On PushT at 4 bits three quantizers tied on the held-out pairs, and only the first in list order was evaluated. E12 (Section~\ref{sec:bits}) replaced this same-bit comparison with the smallest usable bit width.

\begin{table}[h]
\caption{E5: damage of the best post-training quantizer $B$ with \dtm and with \shiter, and the paired difference (TD-MPC2 90 pairs, points or \% of return; DINO-WM 100 pairs, points). \checkmark: lower bound above zero (H3b).}
\label{tab:h3b}
\centering\scriptsize
\setlength{\tabcolsep}{4pt}
\begin{tabular}{@{}lrlrrlc@{}}
\toprule
Task & $k$ & $B$ & $D_{B+\dtm}$ & $D_{B+\shiter}$ & $D_{B+\dtm}-D_{B+\shiter}$ & \\
\midrule
\inputrows{tables/tabG_h3b.tex}
\bottomrule
\end{tabular}
\end{table}

\textbf{Bias correction.} The variant \detbc subtracts the per-unit mean of \dtm's error, estimated on held-out rollouts (Appendix~\ref{app:details}); it is the standard bias correction of post-training quantization. As a remedy it is partial: it repairs \nv{e11.mw.nfix.detbc} of \nv{e11.mw.ndam} damaged Meta-World cells and \nv{e11.dmc.nfix.detbc} of \nv{e11.dmc.ndam} DMC cells at 5M and \nv{e14.nfix.detbc} of \nv{e14.ndam} tasks at 317M, repairs Wall but is worse than \dtm on PushT. Three properties of the fitted vector qualify these numbers. First, it is fitted in sample: at 317M, a vector fitted on half of the held-out episodes leaves a mean residual of \nv{g.bias.mt.oos}~$\Delta$ in the second dynamics layer on the other half, against \nv{g.bias.mt.oos0}~$\Delta$ without correction, so it removes most but not all of the mean bias out of sample. Second, a zero mean residual does not imply a smaller error: on drawer-open at 317M the second dynamics layer's mean squared error grew from \nv{g.bias.mt.drawer}~$\Delta^2$, although \detbc removed the whole damage there. Third, on PushT the vector does not fit the planning states: it raises the planning-time latent error (one step \nv{d4.err1.det} to \nv{d4.err1.detbc}, five steps \nv{d4.err5.det} to \nv{d4.err5.detbc}) and leaves the propagated five-step error of its own calibration unchanged (\nv{d4.b.rel5}), and both subtracting and adding it (\shiterb) are as harmful as \dtm. A better-fitting bias vector might repair more cells; our evidence for a mean-bias component therefore rests on the unbiased variants, not on \detbc (Section~\ref{sec:components}).

\textbf{Quantization-aware training (QAT; a weak baseline).} The planned TD-MPC2 QAT fine-tuned the dynamics and reward networks of each of 11 primary cells with \dtm fake quantization and a straight-through gradient~\citep{bengio2013ste} (5{,}000 steps, batch 256, learning rate $10^{-4}$, data from 20 unquantized episodes). It never removed the damage: QAT+\dtm lost \nv{g.qat.td.det} points or \% of return, with upper bounds far above 5, and it was worse than the best post-training quantizer in \nv{g.qat.td.worseB} of \nv{g.qat.td.of} cells. A post hoc check showed why: the QAT models were damaged even without quantization, by \nv{g.qat.td.none} points or \% in \nv{g.qat.td.none.n} of \nv{g.qat.td.of} cells, so this baseline is weak and we draw no conclusion about training. Within it, \shiter beat \dtm with a lower bound above zero in \nv{g.qat.td.shiter.better} of \nv{g.qat.td.of} cells. The DINO-WM QAT (PushT at 4 bits, stopped at its time budget) behaved differently: QAT+\dtm lost \nv{g.qat.dw.det} points, \nv{g.qat.dw.vs.e1} points less than \dtm without training, and \shiter on the QAT model lost \nv{g.qat.dw.shiter} points more than \dtm. Training may have adapted the model to deterministic rounding; we did not test this. QAT with dither in the forward pass, or longer training, was not run.

\textbf{Weight quantization (exploratory).} Weights stay in floating point in the main experiments. In an exploratory sweep with activations in FP32, 8-bit weights cause no damage (\nv{g.w.td.w8} points or \% of return), while 4-bit round-to-nearest (RTN) weights alone lose \nv{e5.weights.w4.pts} (four TD-MPC2 tasks). Stochastic rounding (SR) of the 4-bit weights, with $M$ copies cycled over the planner iterations, reduced the damage as $M$ grew from 1 to 4 in all four tasks, and with four copies it beat round-to-nearest on push, pick-place and cheetah-run but not on walker ($D_{\mathrm{RTN}}-D_{\mathrm{SR}}$ of \nv{g.w.td.sr4} points or \%). On DINO-WM PushT at 4 bits, 8-bit weights lost \nv{g.w.dw.w8} points and 4-bit round-to-nearest weights \nv{g.w.dw.w4}, an AWQ-like scheme \nv{g.w.dw.awq}, and stochastic rounding with one, two and four copies \nv{g.w.dw.sr}, worse than round-to-nearest. The directions differ between the families, so we do not extend the activation results to weights.
